%Chargement des paquets
\usepackage{amsmath}
\usepackage{amsfonts}
\usepackage{amssymb}
\usepackage{enumerate}
\usepackage{mathtools}
\usepackage{bbm}
\usepackage{xparse, etoolbox}
\usepackage{enumerate}
\usepackage{enumitem}
\usepackage{mathabx}
\usepackage{babel}[french]

%environment
\usepackage{amsthm}
\usepackage{keytheorems}
\usepackage{hyperref} 

%Interface théorème
\newkeytheorem{lem}[
	name = Lemme
]

\newkeytheorem{prop}[
	name = Proposition
]

\newkeytheorem{thm}[
	name = Théorème
]

\newkeytheorem{coro}[
	name = Corollaire
]

\renewcommand*{\proofname}{Démonstration}

\theoremstyle{definition}
\newtheorem{defn}{Définition}

\theoremstyle{definition}
\newtheorem*{nota}{Notation}

\theoremstyle{definition}
\newtheorem*{limi}{Liminaire}

\theoremstyle{remark}
\newtheorem{rmq}{Remarque}

\theoremstyle{plain}
\newtheorem*{fait}{Fait}


%Ensembles de nombres
\newcommand{\N} {\mathbb{N}}
\newcommand{\Ne}{\N^\ast}
\newcommand{\Z} {\mathbb{Z}}
\newcommand{\Q} {\mathbb{Q}}
\newcommand{\R} {\mathbb{R}}
\newcommand{\Rb}{\overline{\mathbb{R}}}
\newcommand{\Rp}{\R_+}
\newcommand{\Rm}{\R_-}
\newcommand{\K} {\mathbb{K}}
\newcommand{\Ket} {\mathbb{K^{\ast}}}
\newcommand{\Cx}{\mathbb{C}}

%Ensembles, tribus
\newcommand{\A}{\mathbb{A}}
\renewcommand{\P}{\mathcal{P}}
\newcommand{\T}{\mathcal{T}}
\newcommand{\F}{\mathcal{F}}
\newcommand{\C} {\mathcal{C}}
\newcommand{\ouv}{\mathcal{O}}
\newcommand{\B}{\mathcal{B}}
\newcommand{\M}{\mathcal{M}}
\newcommand{\etag}{\mathcal{E}}
\newcommand{\etagOp}{\etag^{+}(\Omega)}
\newcommand{\etagO}{\etag(\Omega)}
%%Lp avant quotientage
\newcommand{\Lic}{\mathcal{L}^1}
\newcommand{\Lpc}{\mathcal{L}^p}
\newcommand{\Linfc}{\mathcal{L}^{\infty}}
\newcommand{\Lifull}{\Lic(\Omega, \B, m)}
\newcommand{\Lpfull}{\Lpc(\Omega, \B, m)}
\newcommand{\Linffull}{\Linfc(\Omega, \B, m)}
%%Lp après quotientage
\newcommand{\Li}{L^1}
\newcommand{\Ld}{L^2}
\newcommand{\Lp}{L^p}
\newcommand{\Linf}{L^{\infty}}
\newcommand{\Listd}{\Li(\Omega, m)} 
\newcommand{\Ldstd}{\Ld(\Omega, m)} 
\newcommand{\Lpstd}{\Lp(\Omega, m)} 

%Quelques opérateurs
\DeclareMathOperator{\codim}{codim}
\DeclareMathOperator{\rg}{rg}
\DeclareMathOperator{\tr}{tr}
\DeclareMathOperator{\vect}{Vect}
\DeclareMathOperator{\ima}{Im}
\DeclareMathOperator{\id}{Id}
\DeclareMathOperator{\sign}{sign}
\DeclareMathOperator{\ext}{ext}
\newcommand{\ind}{\mathbbm{1}}
\newcommand*{\spr}[2]{\left(\,#1\,\middle|\,#2\,\right)}
\newcommand{\interior}[1]{%
  {\kern0pt#1}^{\mathrm{o}}%
}
\DeclareMathOperator{\card}{card}
\DeclareMathOperator{\ppcm}{ppcm}

%Commandes de pure flemme
\newcommand{\veps}{\varepsilon}
\newcommand{\intab}{\int_{a}^{b}}
\newcommand{\intR}{\int_{-\infty}^{\infty}}
\newcommand{\intinf}{\int_{- \infty}^{+ \infty}}
\newcommand{\equi}{\Leftrightarrow}
\newcommand{\Kev}{$\K$-espace vectoriel }
\newcommand{\Kevs}{$\K$-espaces vectoriels }
\newcommand{\Rev}{$\R$-espace vectoriel }
\newcommand{\Revs}{$\R$-espaces vectoriels }
\newcommand{\Cev}{$\Cx$-espace vectoriel }
\newcommand{\Cevs}{$\Cx$-espaces vectoriels }
\newcommand{\evn}{(E, \norm{.})}
\newcommand{\interf}[2]{[#1,#2]}
\newcommand{\limb}{\lim\limits_}
\newcommand{\lebn}{\lambda_n}
\newcommand{\pp}{\text{~presque partout}}
\newcommand{\derivp}[2]{\frac{\partial #1}{\partial #2}}
\newcommand{\famiu}{(u_i)_{i \in I}}
\newcommand{\sev}{\text{~s.e.v.}}
\newcommand{\Mnp}{M_{n,p}(\K)}
\newcommand{\Mn}{M_{n}(\K)}
\newcommand{\tq}{\text{~t.q.~}}
\newcommand{\mq}{~montrer~que~}
\newcommand{\et}{\quad \text{et} \quad}
\newcommand{\ou}{\text{~ou~}}
\newcommand{\Kx}{\K[X]}


%Commandes de gars chelou
\renewcommand{\subset}{\subseteq}

%Commande restriction, ty duckduckgo
\def\restr#1#2{\mathchoice
              {\setbox1\hbox{${\displaystyle #1}_{\scriptstyle #2}$}
              \restraux{#1}{#2}}
              {\setbox1\hbox{${\textstyle #1}_{\scriptstyle #2}$}
              \restraux{#1}{#2}}
              {\setbox1\hbox{${\scriptstyle #1}_{\scriptscriptstyle #2}$}
              \restraux{#1}{#2}}
              {\setbox1\hbox{${\scriptscriptstyle #1}_{\scriptscriptstyle #2}$}
              \restraux{#1}{#2}}}
\def\restraux#1#2{{#1\,\smash{\vrule height .8\ht1 depth .85\dp1}}_{\,#2}} 

%Quelques normes
\DeclarePairedDelimiter\abs{\lvert}{\rvert}
\DeclarePairedDelimiter\labs{\bigl\lvert}{\bigr\rvert}
\DeclarePairedDelimiter\norm{\lVert}{\rVert}
\DeclarePairedDelimiterXPP\normi[1]{}\lVert\rVert{_1}{\ifblank{#1}{\:\cdot\:}{#1}}
\DeclarePairedDelimiterXPP\normd[1]{}\lVert\rVert{_2}{\ifblank{#1}{\:\cdot\:}{#1}}
\DeclarePairedDelimiterXPP\normp[1]{}\lVert\rVert{_p}{\ifblank{#1}{\:\cdot\:}{#1}}
\DeclarePairedDelimiterXPP\normq[1]{}\lVert\rVert{_q}{\ifblank{#1}{\:\cdot\:}{#1}}
\DeclarePairedDelimiterXPP\norminf[1]{}\lVert\rVert{_\infty}{\ifblank{#1}{\:\cdot\:}{#1}}

%Bracket en angle
\DeclarePairedDelimiter\angl{\langle}{\rangle}

%Set, parenthèses
\newcommand{\set}[1]{\{ #1 \}}
\newcommand{\bigset}[1]{\bigl\{ #1 \bigr\}}
\newcommand{\Bigset}[1]{\Bigl\{ #1 \Bigr\}}
\newcommand{\bigpar}[1]{\bigl( #1 \bigr)}
\newcommand{\Bigpar}[1]{\Bigl( #1 \Bigr)}

%Opitmisation des opérateurs de comparaison
\DeclarePairedDelimiter\tio{o (}{)}
\DeclarePairedDelimiter\gro{O (}{)}


\pagestyle{fancy}

\fancyhead[C]{}
\fancyhead[R]{}

\begin{document}

\underline{Source :} Tauvel - Analyse complexe - pages 126-128 puis 135

\underline{Leçons :}
\begin{itemize}
\item 
\end{itemize}

\begin{nota}
On note $U$ un ouvert de $\Cx$.
\end{nota}

\begin{defn}[arc, lacet]
Un arc $\gamma$ de $U$ est une application continue d'un compact de $[x,y] \subset \R$ dans $U$.
L'origine d'un arc est l'image de $x$ par $\gamma$ et son extrémité l'image de $y$.
Un arc est appelé lacet si origine et extrémité sont confondus.
\end{defn}

\begin{rmq}
Par changement de variable $t' = (1-t)x + ty$ on peut toujours se ramener au compact $[0,1]$. 
Dans la suite, on considère donc tous les arcs définis sur [0,1].
\end{rmq}

\begin{defn}[déformation]
Soient $\gamma_1, \gamma_2$ deux arcs de $U$. On appelle déformation de $\gamma_1$ à $\gamma_2$ toute application continue :
\[ \begin{array}{c c c c}
\delta : & [0,1]^2 & \to & U \\
& (t,u) & \mapsto & \delta(t,u) 
\end{array} \]
vérifiant $\delta(t,0) = \gamma_1(t)$ et $\delta(t,1) = \gamma_2(t)$.
\end{defn}

\begin{defn}[homotopie]
\begin{enumerate}[label=(\roman*)]
\item Deux arcs $\gamma_1, \gamma_2$ de $U$ ayant même origine $a$ et même extrémité $b$ sont dits homotopes dans $U$
s'il existe une déformation $\delta$ de $\gamma_1$ à $\gamma_2$ telle que :
\[ \forall u \in [0,1], \quad \delta(0,u)=a \et \delta(1,u)=b . \]
S'il en est ainsi on dit que $\delta$ est une homotopie de $\gamma_1$ à $\gamma_2$.
\item Deux lacets $\gamma_1, \gamma_2$ de $U$ sont dits homotopes dans $U$ s'il existe une déformation $\delta$ 
de $\gamma_1$ à $\gamma_2$ telle que :
\[ forall u \in [0,1], \quad \delta(0,u) = (1,u) . \]
S'il en est ainsi on dit que $\delta$ est une homotopie de $\gamma_1$ à $\gamma_2$.
\end{enumerate}
\end{defn}

\begin{lem}
\label{lem:1}
Soient $\gamma_1, \gamma_2$ des arcs (resp. des lacets) de $U$ homotopes dans $U$.
Alors $\gamma_1$ et $\gamma_2$ appartiennent à la même composante connexe de $U$.
\end{lem}

\begin{proof}
Soit $\delta$ la déformation permettant de passer d'un lacet à l'autre et soit $t \in [0,1]$.
Comme [0,1] est connexe et $\delta(t,u)$ continue en $u$, l'ensemble $\delta(t, [0,1])$ est un connexe de $U$.
Donc $\gamma_1(t)$ et $\gamma_2(t)$ appartiennent à la même classe d'équivalence et ce résultat est vrai pour tout $t \in [0,1]$.
\end{proof}

\begin{thm}[Cauchy pour un convexe, admis]

\begin{thm}[Cauchy]
Soient $\gamma_1, \gamma_2$ des arcs (resp. des lacets) de $U$ homotopes dans $U$.
Pour toute fonction $f$ holomorphe dans $U$, on a :
\[ \int_{\gamma_1} f(z) dz = \int_{\gamma_2} f(z) dz . \]
\end{thm}

\begin{proof}
Selon le lemme $\ref{lem:1}$, les deux arcs (reps. lacets) appartiennent à la même composante connexe de $U$, 
on peut donc toujours supposer $U$ connexe.
Si $U = \Cx$, on prend $\veps >0$ quelconque, sinon on prend $\veps = d(\im{\gamma_1}, \Cx \setminus U)$.
Dans les deux cas, le disque $D_t = D(\gamma_1(t), \veps)$ est contenu dans $U$.
Par continuité de $\gamma_1$, il existe $\eta > 0$ tel que :
\[ \forall t, t' \in [0,1], \abs{t-t'} < \eta \Rightarrow \abs{\gamma_1(t) - \gamma_1(t')} < \veps . \]
Posons $0=t_0 < t_1 < \dots < t_n = 1$ une subdivision de [0,1] de pas maximal strictement inférieur à $\eta$.

\end{proof}

\end{document}